% Generated draft: requires mathematical and source review.
\documentclass[11pt]{article}
\usepackage[margin=25mm]{geometry}
\usepackage{fontspec}
\setmainfont{Latin Modern Roman}
\usepackage{amsmath,amssymb,mathtools,unicode-math}
\setmathfont{Latin Modern Math}
\usepackage{enumitem,needspace,xurl,hyperref}
\setlength{\parindent}{0pt}
\setlength{\parskip}{5pt}
\setlength{\emergencystretch}{3em}
\newcommand{\N}{\mathbb{N}}
\newcommand{\Z}{\mathbb{Z}}
\newcommand{\Q}{\mathbb{Q}}
\newcommand{\R}{\mathbb{R}}
\newcommand{\C}{\mathbb{C}}
\newcommand{\F}{\mathbb{F}}
\newcommand{\A}{\mathbb{A}}
\newcommand{\PP}{\mathbb{P}}
\newcommand{\E}{\mathbb{E}}
\newcommand{\eps}{\varepsilon}
\newcommand{\dd}{\,\mathrm{d}}
\begin{document}
\section*{UnsolvedMath ID: 10400153}
\textbf{Source problem code:} catalogue entry
\subsection{Definitions and mathematical statement}
For an abelian group \(H\), an Eilenberg--Mac Lane space \(K(H,2)\) has \(\pi_2\cong H\) and all other positive homotopy groups zero. Maps from suitable manifold/CW inputs to this target are classified by degree-two cohomology classes with coefficients in \(H\).

An HQFT is a homotopy quantum field theory: in a standard functorial formulation it assigns state spaces to structured manifolds equipped with maps to a target and linear maps to structured cobordisms carrying compatible target maps, with gluing, tensor and homotopy-invariance axioms. For spin or \(\operatorname{spin}^c\) HQFTs the cobordisms and boundary objects have the corresponding tangential structures.

Study these theories with target \(K(H,2)\) in all three spacetime dimensions
\[
 1+1,\qquad2+1,\qquad3+1,
\]
for the group rendered as ``\(\mathbb Z N\)'' in the extraction. If this means \(H=\mathbb Z/N\mathbb Z\), it is the cyclic-group version; \(\mathbb Z^N\) would instead be a free-abelian-group target. The missing subscript/superscript typography must be confirmed, and neither reading is silently chosen.

The request ``study'' does not prescribe a construction, classification equivalence, output invariants, coefficient category, unitarity or extendedness convention. It retains both spin and \(\operatorname{spin}^c\) variants and each dimension; no known lower-dimensional or aspherical-target theory is substituted for the \(K(H,2)\) problem. The summary provides only an open/triage label and no precise theorem, so current status requires primary-literature review.
\subsection{Short English statement}
Develop spin and \(\operatorname{spin}^c\) HQFTs with target \(K(H,2)\) in dimensions two, three and four. The group typography and the intended construction/classification scope are missing, and all three dimensions remain separate targets.
\subsection{Sources}
\begin{itemize}
\item Local registry: \url{data/unsolvedmath/problems.json}, permanent UnsolvedMath ID \(10400153\), source code catalogue entry; supplied statement, background, and research summary, updated 2026-08-06.
\item Ohtsuki, \emph{Problems on invariants of knots and 3-manifolds} (2002), Problem 8.7, attributed to V. Turaev, PDF page 124: \url{https://msp.org/gtm/2002/04/gtm-2002-04-024s.pdf}. Supplied PDF-font-structure extraction accessed 2026-07-29; provenance distinguishes 231 numbered items.
\item No exact group-typography recovery, detailed HQFT convention or separate current theorem reference is supplied.
\end{itemize}
\end{document}
